\section{Introduction}
\footnote{GitHub link: \url{http://}}

Every year, thousands of machine learning models are built to predict clinical outcomes  from heart disease risk to cancer prognosis to treatment response (Collins et al., 2024). These models use a range of techniques, from straightforward logistic regression to complex deep learning, drawing on electronic health records, genomic data, imaging, and clinical trials (Greener et al., 2022). The way these models are usually tested is by holding back part of the training data and checking how well the model predicts on that held-out portion. The problem is that this only tells you how the model performs on data that looks like what it was trained on. It says nothing about what will happen when the model meets patients from a different hospital, a different age group, or a different country.
The evidence shows that what usually happens is the model fails. Chekroud et al. (2024) built models to predict antipsychotic treatment outcomes using data from five independent clinical trials. When tested within each trial, the models scored 0.72 balanced accuracy. With proper cross-validation, this dropped to 0.60. When tested on a completely separate trial, it fell to 0.54  no better than guessing. The models had not learned anything about treatment response; they had picked up on patterns specific to each trial that did not carry over to new data. This kind of failure is not limited to psychiatry. The same pattern has been found in genomics (Whalen et al., 2022), clinical biomarker research (Wörheide et al., 2021), immunology (Springer & Doering, 2021), and molecular prediction (Ektefaie et al., 2024). The reasons are well understood. Some theoretical fixes have been proposed. But what nobody has built yet is a practical tool that can tell you, before you deploy a model, whether it is likely to work in a specific new setting. That is the gap this portfolio addresses.
